\section{The first disk near a hexagonal front}\label{sec:static}

A rhombic lattice front whose notch hinge is exactly $60^\circ$ is a piece of the hexagonal packing:
every disk of the front touches its neighbours in the triangular lattice. This section shows that such
fronts are unstable in the direction the conjecture cares about. Perturb a few hinges by amounts of
order $t$ and let the disks shrink at a rate of order $t$, and the very first disk can be a pentagon.

\subsection{The notch-60 lattice}
\begin{lemma}[Notch-60 law]\label{lem:notch60}
A rhombic lattice front with $m$ down vectors and $n$ up vectors of bearing $u$, the down vectors at
$u-120^\circ$ (notch hinge $60^\circ$), closes if and only if
\[
\tan u=\frac{\sqrt3\,m}{2n-m}.
\]
Both families point to the right, $30^\circ<u<90^\circ$, exactly when $n/2<m<2n$.
\end{lemma}

\begin{proof}
Expand $\sin(u-2\pi/3)=-\tfrac12\sin u-\tfrac{\sqrt3}2\cos u$: the closing condition
$m\sin(u-2\pi/3)+n\sin u=0$ becomes $(2n-m)\sin u=\sqrt3\,m\cos u$. The bounds on $u$ are $\tan u>1/\sqrt3$
and $2n-m>0$.
\end{proof}

\noindent For $(2,3)$ this gives $\tan u=\sqrt3/2$, $u\approx40.89^\circ$ and down bearing
$\approx-79.11^\circ$; for $(1,2)$ it gives a vertical down vector, which is why the construction below
needs $m,n\ge2$.

\begin{definition}\label{def:hexfront}
Let $m,n\ge2$ with $n/2<m<2n$. Write the front word as $D_1\cdots D_m\,U_1\cdots U_n$ and name the disks
around the notch $Z\xrightarrow{D_{m-1}}A\xrightarrow{D_m}B\xrightarrow{U_1}C\xrightarrow{U_2}E$. The front
$H(\alpha,\sigma,\eps,c)$ has all its disks, the seeds included, of diameter~\eqref{eq:diam}, consecutive disks
touching, and the hinges of the notch-60 lattice except four:
\begin{gather*}
h(D_{m-1},D_m)=180^\circ+\alpha,\qquad h(D_m,U_1)=60^\circ+\sigma,\qquad h(U_1,U_2)=180^\circ-\eps,\\
h(U_n,D_1)=300^\circ-(\alpha+\sigma-\eps).
\end{gather*}
\end{definition}

The last hinge, at the ridge, is forced: the deviations $h-180^\circ$ of the hinges of a closed front sum
to zero. To fix scale and vertical translation, set the diameter of the first disk to $1$ at height $0$.
The bearings have the form $\phi+\beta_j(\alpha,\sigma,\eps)$, where the hinges determine the relative
bearings $\beta_j$. Generate the edge lengths by~\eqref{eq:edge-exact}. At zero parameters the
vertical closing error has derivative $\sum_j\cos\beta_j=\Lambda_0>0$ with respect to $\phi$.
The implicit function theorem therefore determines an analytic closing rotation for small
parameters. Horizontal closure defines the positive circumference. The last tangency is included
in this recurrence, since its endpoint returns to height $0$ and hence to diameter $1$.
Both $\eta_U$ and $\eta_D$ are $O(|\alpha|+|\eps|+|c|)$, so their maximum has this bound and
their product has its square as a bound. The one notch angle is
$60^\circ+\sigma$, and the rate is $|c|$.

At $t=0$ the disks of $H$ sit on the triangular lattice $\{xD+yU:x,y\in\mathbb Z\}$, where $D$ and $U$ now
denote the unit down and up vectors. Since $D\cdot U=-\tfrac12$,
\begin{equation}\label{eq:eisenstein}
|xD+yU|^2=x^2-xy+y^2,
\end{equation}
which takes the values $0,1,3,4,7,\dots$ and never $2$. Two disks of $H$ that are not neighbours are at
distance at least $\sqrt3$ at $t=0$, except two pairs at distance $1$: the disks on either side of the
notch, $A$ and $C$, and the disks on either side of the ridge. Indeed, going right from one disk to the
other along the front adds the vectors of a proper cyclic subword of $D^mU^n$, with sum $xD+yU$,
$0\le x\le m$, $0\le y\le n$. By~\eqref{eq:eisenstein} this has length $1$ only for $(x,y)$ equal to
$(1,0)$ or $(0,1)$, which are neighbours, or $(1,1)$, which is the subword $DU$ at the notch or $UD$ at
the ridge. The other images of the second disk add $k(mD+nU)$. For $k=-1$ the result is minus the sum
of the complementary subword, covered by the same list. For $k\ge1$ or $k\le-2$ both coordinates are
at least $2$ or both at most $-2$, and since $x^2-xy+y^2=\tfrac14(x+y)^2+\tfrac34(x-y)^2$ the squared
length is at least $4$.

\begin{lemma}[Validity]\label{lem:valid}
Along $(\alpha,\sigma,\eps,c)=t(a,s,e,-q)$ the clearances of these two pairs are, to first order,
$\tfrac{\sqrt3}2\bigl(s+\tfrac q2\kappa_u\bigr)t$ at the notch and $\tfrac{\sqrt3}2\bigl(a+s-e-\tfrac q2\kappa_u\bigr)t$ at the
ridge, with $\kappa_u=\cos(u-60^\circ)$. So the lower bounds in~\eqref{eq:window} make $H(ta,ts,te,-tq)$ a
valid front for small $t>0$.
\end{lemma}

\begin{proof}
Let $T$ be the middle disk of either triple and $L$, $R$ its neighbours. Write $\ell_1=|T-L|$,
$\ell_2=|T-R|$, and let $\phi$ be the angle of the triangle $LTR$ at $T$: the angle above the chain at the
notch and below it at the ridge. Then $\phi$ is $60^\circ$ plus a hinge deviation, which the
definition prescribes exactly: $\phi=60^\circ+\sigma$ at the notch and $\phi=60^\circ+(\alpha+\sigma-\eps)$ at the
ridge. By the law of cosines,
\[
|L-R|^2=\ell_1^2+\ell_2^2-2\ell_1\ell_2\cos\phi .
\]
At $t=0$, $\ell_1=\ell_2=|L-R|=1$ and $\phi=60^\circ$, so
$\tfrac{d}{dt}|L-R|=\tfrac12(\dot\ell_1+\dot\ell_2)+\tfrac{\sqrt3}2\dot\phi$. Tangency gives
$\ell_1=\tfrac12(d_T+d_L)$ and $\ell_2=\tfrac12(d_T+d_R)$ exactly. The clearance of the pair is
$g=|L-R|-\tfrac12(d_L+d_R)$, so
\[
\dot g=\tfrac{\sqrt3}2\dot\phi+\tfrac12\Bigl(\dot d_T-\tfrac12(\dot d_L+\dot d_R)\Bigr).
\]
With $c=-qt$ and heights $y_k$ at $t=0$, $\dot d_k=-q\,(y_k-y_0)$, and the reference height $y_0$ cancels.
At the ridge $T$ is reached by $U_n$ and left by $D_1$, so
$y_T-\tfrac12(y_L+y_R)=\tfrac12(\Imag U-\Imag D)=\tfrac12\bigl(\sin u-\sin(u-120^\circ)\bigr)=\tfrac{\sqrt3}2\kappa_u$; at the
notch the same expression is $-\tfrac{\sqrt3}2\kappa_u$. Hence
\[
\dot g_{\mathrm{notch}}=\tfrac{\sqrt3}2\bigl(s+\tfrac q2\kappa_u\bigr),\qquad
\dot g_{\mathrm{ridge}}=\tfrac{\sqrt3}2\bigl(a+s-e-\tfrac q2\kappa_u\bigr).
\]
Every other pair of non-neighbours is at distance at least $\sqrt3$ at $t=0$, by the remark before the
lemma, so it stays clear for small $t$. Neighbours touch by construction. Under the lower bounds
in~\eqref{eq:window} both derivatives are positive, so $g>0$ at both pairs for small $t>0$.
\end{proof}

\noindent Shrinking closes the ridge: the hexagonal front itself is not a valid front once $c<0$, and the
ridge has to open by at least $\tfrac q2\kappa_u t$. With the other hinges at their lattice values this forces
$\alpha+\sigma-\eps>0$, so with $\sigma<\eps$ the down vectors must be irregular too ($a>0$). On fronts with
$m+n\ge5$ one can instead perturb another hinge inside a run, which changes nothing at $Q_0$.

Theorem~\ref{thm:static} of the introduction says that $H(ta,ts,te,-tq)$ places a pentagon first when
$(a,s,e,q)$ satisfies~\eqref{eq:window}. The rest of this section proves it.

\subsection{The lowest site}
\begin{figure}[t]
\centering
\begin{tikzpicture}[scale=1.55,every node/.style={font=\small}]
\coordinate (Z) at (1,1.7321); \coordinate (A) at (0.5,0.8660); \coordinate (B) at (0,0);
\coordinate (C) at (1,0); \coordinate (E) at (2,0); \coordinate (Q) at (1.5,0.8660);
\foreach \p in {Z,A,B,C,E} \draw[fill=black!8] (\p) circle (0.5);
\draw[thick] (Q) circle (0.5);
\foreach \p in {Z,A,B,C,E,Q} \fill (\p) circle (0.9pt);
\draw[-{Latex[length=2mm]}] (Z)--(A); \draw[-{Latex[length=2mm]}] (A)--(B);
\draw[-{Latex[length=2mm]}] (B)--(C); \draw[-{Latex[length=2mm]}] (C)--(E);
\draw[densely dashed] (Z)--(E);
\node[above left=7pt] at (Z) {$Z$}; \node[left=9pt] at (A) {$A$}; \node[below=9pt] at (B) {$B$};
\node[below=9pt] at (C) {$C$}; \node[below=9pt] at (E) {$E$}; \node at ($(Q)+(0.17,-0.2)$) {$Q_0$};
\draw (0.3,0) arc (0:60:0.3); \node at (0.42,0.2) {$\tau$};
\draw (1.25,0) arc (0:180:0.25); \node at (1,0.36) {$\gamma$};
\node[align=left,anchor=west] at (2.75,1.1) {$\tau=60^\circ+\sigma$\\$\gamma=180^\circ-\eps$\\$h_A=180^\circ+\alpha$};
\end{tikzpicture}
\caption{The notch of a hexagonal front, turned so that $U_1=C-B$ is horizontal. The site $Q_0$ touches
four front disks, $Z$, $A$, $C$ and $E$, and $Z$, $E$ are two diameters apart (dashed). The pentagon disk
touches $A$ and $E$; at constant diameter it must clear $C$, which needs $\eps>\sigma$, and fit between $Z$
and $E$, which needs $2\sigma+\alpha>\eps$.}\label{fig:hex}
\end{figure}

\begin{lemma}\label{lem:hexsite}
At $\alpha=\sigma=\eps=c=0$, with unit diameters:
\begin{enumerate}
\item[(a)] every position touching two front disks is a lattice point;
\item[(b)] among lattice points strictly above the front, the lowest are $Q_0=A+U$ and its periodic
images, and every other one is higher by at least $\min(|\Imag D|,\Imag U)$;
\item[(c)] $Q_0$ touches exactly the four front disks $Z,A,C,E$; $|Z-E|=2$ with $Q_0$ the midpoint, and
$|Q_0-B|=\sqrt3$.
\end{enumerate}
\end{lemma}

\begin{proof}
(a) A position touching disks at $P$ and $R$ has $|Q-P|=|Q-R|=1$, so $|P-R|\le2$, and
by~\eqref{eq:eisenstein} $|P-R|$ is $1$, $\sqrt3$ or $2$. The lattice is invariant under rotation by
$60^\circ$. In the first case $Q=P+(R-P)e^{\pm i\pi/3}$. In the second, the six solutions
of $x^2-xy+y^2=3$ show that $R-P=w+we^{i\pi/3}$ for a unit lattice vector $w$, and the two positions are
$P+w$ and $P+we^{i\pi/3}$. In the third, the solutions of $x^2-xy+y^2=4$ are twice the unit vectors, and
$Q$ is the midpoint. In each case $Q$ is a lattice point.

(b) Use coordinates $xD+yU$ and put $\delta=-\Imag D>0$, $\upsilon=\Imag U>0$. The height is
$-x\delta+y\upsilon$, and closure of the period $mD+nU$ gives $m\delta=n\upsilon$. Take $B$ as the origin and
$A=(-1,0)$, so the front vertices are $(x,0)$ for $-m\le x\le0$ and $(0,y)$ for $0\le y\le n$, repeated
with period $(m,n)$. The map $(x,y)\mapsto xD+yU$ preserves orientation and $\Real D,\Real U>0$, so the front is a
graph over the horizontal axis, and a point lies above it exactly when it lies to the left of the
staircase traversed in the direction of increasing $x$ and $y$. Write a lattice point as
$(km-j,\,y)$ with $k\in\mathbb Z$ and $1\le j\le m$. It is strictly above the front exactly when
$y\ge kn+1$. Relative to $B$, and using $km\delta=kn\upsilon$, its height is
\[
-(km-j)\delta+y\upsilon=j\delta+(y-kn)\upsilon\;\ge\;j\delta+\upsilon .
\]
This is smallest, $\delta+\upsilon$, exactly for $j=1$, $y=kn+1$, which is
$(-1,1)+k(m,n)=Q_0+k(mD+nU)$. Any other point has $j\ge2$ or $y\ge kn+2$, and so is higher by at least
$\delta$ or $\upsilon$.

(c) The six unit neighbours of $Q_0=(-1,1)$ are $Q_0+(1,0)=(0,1)=C$, $Q_0-(0,1)=A$,
$Q_0+(1,1)=(0,2)=E$, $Q_0-(1,1)=(-2,0)=Z$, and $(-1,2)$, $(-2,1)$. The last two are not front vertices,
because $m,n\ge2$. Then $Z-E=-2(D+U)$ has length $2$ with midpoint $Q_0$, and $Q_0-B=U-D$ has length
$\sqrt3$ by~\eqref{eq:eisenstein}.
\end{proof}

\noindent So six pairs share the apex $Q_0$: $ZA$ (a triangle on $D_{m-1}$), $AC$ (the rhombus), $CE$ (a
triangle on $U_2$), $AE$ and $ZC$ (pentagons) and $ZE$ (a transition removing three vectors). The first
disk is decided among them at first order. The pair $ZE$ is degenerate and is treated by
its separation, rather than by a transverse-apex expansion.

\subsection{First-order clearances}
Turn and translate the configuration so that $B=0$ and $U_1$ points along $1$; clearances do not change,
and the height enters only multiplied by $c$, so its first-order change does not matter. Then
$A=\omega$, $C=1$, $E=2$, $Z=2\omega$, $Q_0=1+\omega$ with $\omega=e^{i\pi/3}$, and the contact normals
$n_k=Q_0-z_k$ are
\[
n_A=1,\qquad n_C=\omega,\qquad n_E=\omega-1,\qquad n_Z=1-\omega.
\]
Along $(\alpha,\sigma,\eps,c)=t(a,s,e,-q)$ write $\dot z$ for $dz/dt$ at $t=0$. For a candidate family $Q(t)=Q_0+t x+o(t)$, differentiating the clearance $|Q-z_k|-\tfrac12(d_k+d_Q)$ as in~\eqref{eq:tangency} shows that, to
first order, the disk is clear of $z_k$ to strict first order when $n_k\cdot x>\lambda_k$; a tangency requires $n_k\cdot x=\lambda_k$, with
\[
\lambda_k=n_k\cdot\dot z_k-\tfrac q2\,h_k
\]
and $h_k$ the height of $z_k$ above $Q_0$ (diameters normalised to $1$ at the height of $Q_0$).

\begin{lemma}\label{lem:hexclear}
Put $\mu_C=\lambda_A+\lambda_E-\lambda_C$ and $\mu_Z=-(\lambda_Z+\lambda_E)$. To first order:
\begin{enumerate}
\item[(a)] the disk touching $A$ and $E$ clears $C$ by $\mu_C$ and $Z$ by $\mu_Z$;
\item[(b)] the disks touching $A,C$ and $C,E$ overlap $E$ and $A$ by $\mu_C$, and the disk touching $Z,C$
overlaps $A$ by $\mu_C+\mu_Z$;
\item[(c)] if $\mu_Z>0$ no disk near $Q_0$ touches both $Z$ and $E$, and the disk touching $Z,A$ is higher than the one
touching $A,E$ by $\tfrac2{\sqrt3}\mu_Z\cos u$;
\item[(d)] $\mu_C=\tfrac{\sqrt3}2(e-s)-\tfrac q8\bigl(3\sin u+5\sqrt3\cos u\bigr)$ and
$\mu_Z=\tfrac{\sqrt3}2(2s+a-e)+\tfrac q2\bigl(3\sin u+\sqrt3\cos u\bigr)$, per unit $t$.
\end{enumerate}
\end{lemma}

\begin{proof}
Parts (a) to (c) use only $n_E=n_C-n_A$ and $n_Z=-n_E$. At the vertex $n_A\cdot x=\lambda_A$, $n_E\cdot x=\lambda_E$ one gets
$n_C\cdot x-\lambda_C=\lambda_A+\lambda_E-\lambda_C$ and $n_Z\cdot x-\lambda_Z=-\lambda_E-\lambda_Z$. At $n_A\cdot x=\lambda_A$, $n_C\cdot x=\lambda_C$ the
clearance to $E$ is $\lambda_C-\lambda_A-\lambda_E=-\mu_C$, and symmetrically for $C,E$ against $A$. At $n_Z\cdot x=\lambda_Z$,
$n_C\cdot x=\lambda_C$ the clearance to $A$ is $\lambda_Z+\lambda_C-\lambda_A=-\mu_C-\mu_Z$.

The pair $ZE$ has opposite normals, so it has no transverse apex and needs a separate argument. Suppose
a disk at $Q$ touches both, with $Q\to Q_0$ as $t\to0$. By the triangle inequality
\[
|Z-E|\le|Q-Z|+|Q-E|=\tfrac12(d_Z+d_E)+d_Q,
\]
so $g(t)=|Z-E|-\tfrac12(d_Z+d_E)-d_Q\le0$. Normalise diameters to $1$ at the height of $Q_0$ at $t=0$;
then $d_Q=1+c\,(\Imag Q-\Imag Q_0)=1+o(t)$. Since $Z-E=-2n_Z$ with $|n_Z|=1$ at $t=0$,
\[
g(t)=t\bigl(-n_Z\cdot\dot Z-n_E\cdot\dot E+\tfrac q2(h_Z+h_E)\bigr)+o(t)=t\,\mu_Z+o(t),
\]
using $n_E=-n_Z$ and $\dot d_k=-qh_k$. So $\mu_Z>0$ gives $g(t)>0$ for small $t>0$, a contradiction. No
apex expansion is used, so this covers candidates whose distance from $Q_0$ is much larger than $t$.
The vertices $ZA$ and $AE$ both lie on
the line $n_A\cdot x=\lambda_A$, a distance $\tfrac2{\sqrt3}\mu_Z$ apart along $in_A$, whose height in the cylinder
frame is $\cos u>0$.

(d) The hinges give $\dot A=is\,\omega$, $\dot C=0$, $\dot E=ie$, $\dot Z=i(2s+a)\omega$ at $q=0$, so
$\lambda_A=-\tfrac{\sqrt3}2s$, $\lambda_C=0$, $\lambda_E=\tfrac{\sqrt3}2e$ and $\lambda_Z=-\tfrac{\sqrt3}2(2s+a)$. The rate stretches
each side to the mean of its two diameters, adding $-q\ell_{BC}$ to $\dot C$, $-q(\ell_{BC}+\ell_{CE})$ to $\dot E$,
$-q\ell_{AB}\omega$ to $\dot A$ and $-q(\ell_{AB}+\ell_{ZA})\omega$ to $\dot Z$, where $\ell_{XY}=\tfrac12(h_X+h_Y)$.
Collecting terms,
\[
\mu_C^{q}=-\tfrac q4\,(3h_A-h_B-5h_C+h_E),\qquad \mu_Z^{q}=-\tfrac q4\,(2h_A+2h_B+2h_C-h_Z-h_E),
\]
and the heights are $h_z=-\Imag\bigl(e^{iu}(Q_0-z)\bigr)$: $h_A=-\sin u$, $h_B=-\sqrt3\sin(u+30^\circ)$,
$h_C=-\sin(u+60^\circ)$, $h_E=-\sin(u+120^\circ)$, $h_Z=-\sin(u-60^\circ)$. Since
$\sin(u-60^\circ)+\sin(u+120^\circ)=0$ this gives (d).
\end{proof}

\begin{proof}[Proof of Theorem~\ref{thm:static}]
By Lemma~\ref{lem:valid} the front is valid for small $t>0$. By~(d) the upper bound in~\eqref{eq:window} is
$\mu_C>0$. Since $\tfrac12(3\sin u+\sqrt3\cos u)=\sqrt3\kappa_u$, (d) and Lemma~\ref{lem:valid} give $\mu_Z$ as the sum
of the notch and ridge clearance rates plus $\sqrt3q\kappa_u$, so $\mu_Z>0$ on every valid front.

\emph{The pentagon disk is a valid position.} The normals $n_A$ and $n_E$ are not parallel, so
Lemma~\ref{lem:apex} gives a $C^1$ family $Q_{AE}(t)\to Q_0$ touching $A$ and $E$. By
Lemma~\ref{lem:hexclear}(a) its clearances from $C$ and $Z$ are $t\mu_C+o(t)$ and $t\mu_Z+o(t)$, both
positive for small $t$. Every other front disk, periodic images included, is at distance at least
$\sqrt3$ from $Q_0$ at $t=0$ (Lemma~\ref{lem:hexsite}(c) and~\eqref{eq:eisenstein}), so it stays clear.
$Q_0$ is strictly above the front, and the front moves continuously with $t$, so $Q_{AE}(t)$ is above it.

\emph{Nothing else is as low.} Suppose that for some sequence $t_k\to0$ there is a valid position
$Q_k\ne Q_{AE}(t_k)$ with $\Imag Q_k\le\Imag Q_{AE}(t_k)$, touching front disks $P_k$ and $R_k$. We derive a
contradiction in three steps.

\smallskip\noindent\emph{Step 1: compactness.} Normalise the diameter to $1$ at height $0$, so the front
disks have diameters in $[\tfrac12,2]$ and heights in a fixed bounded interval for small $t$. If $Q$ touches
$P$, then $|\Imag Q-\Imag P|\le|Q-P|=d_P+\tfrac c2(\Imag Q-\Imag P)$, so $|\Imag Q-\Imag P|\le2d_P\le4$ once $|c|\le\tfrac12$.
Hence $Q_k$ has bounded height, and modulo the period it lies in a bounded set. Its parents lie within
distance $4$ of it, and the circumference is bounded below, so up to periodic translation $P_k$ and
$R_k$ range over a finite set of front disks and images. Passing to a subsequence, the pair
$(P_k,R_k)$ is fixed and $Q_k\to Q^*$.

\smallskip\noindent\emph{Step 2: the limit is $Q_0$.} Touching and non-overlap are closed conditions, and
all diameters tend to $1$. So $Q^*$ is at distance exactly $1$ from the limits of $P_k$ and $R_k$, and at
distance at least $1$ from every front disk at $t=0$. By Lemma~\ref{lem:hexsite}(a), $Q^*$ is a lattice
point. It lies on or above the limiting front. The only lattice points on the front polyline are its
vertices, which are front disks at distance $0$, so $Q^*$ is strictly above. Finally
$\Imag Q^*\le\lim\Imag Q_{AE}(t_k)=\Imag Q_0$. By Lemma~\ref{lem:hexsite}(b), $Q^*$ is $Q_0$ up to periodic
translation, and by Lemma~\ref{lem:hexsite}(c) the limiting parents are two of $Z,A,C,E$.

\smallskip\noindent\emph{Step 3: the six pairs.} For large $k$, $Q_k$ is a position near $Q_0$ touching one
of the six pairs. Five of them have non-parallel normals at $Q_0$. For these, the implicit function
theorem in the proof of Lemma~\ref{lem:apex}, with $t$ in place of $c$, gives a neighbourhood of $Q_0$ in
which the touching position is unique for small $t$ and lies on the $C^1$ apex family, so $Q_k$ lies on
that family.
\begin{itemize}
\item $AE$: then $Q_k=Q_{AE}(t_k)$, which was excluded.
\item $AC$, $CE$ and $ZC$: by Lemma~\ref{lem:hexclear}(b) these families overlap $E$, $A$ and $A$ by
$t\mu_C+o(t)$, $t\mu_C+o(t)$ and $t(\mu_C+\mu_Z)+o(t)$, so they are not valid for large $k$.
\item $ZA$: by Lemma~\ref{lem:hexclear}(c) it is higher than $Q_{AE}$ by $\tfrac2{\sqrt3}\mu_Z\cos u\,t+o(t)>0$.
\end{itemize}
The sixth pair, $ZE$, admits no touching position near $Q_0$, by the separation argument in the proof of
Lemma~\ref{lem:hexclear}. Every case gives a contradiction. So for small $t>0$ the pentagon disk
touching $A$ and $E$ is the unique lowest valid position, and the first transition is a pentagon.
\end{proof}

\noindent Lemma~\ref{lem:hexclear}, the notch part of Lemma~\ref{lem:valid} and the step from validity to
$\mu_Z>0$ are checked in Lean (\lean{Static.lean}): the four tangencies and three hinge rates of the front
determine the first-order clearances of all six candidates and of the notch pair. The file
\lean{StaticExists.lean} constructs the local five-disk and apex families for the prescribed
rates. Periodic closure and global exclusion are supplied by the analytic arguments above.

On valid rays above the window ($\mu_C<0$), the tested fronts place a rhombus $AC$. The independent
simulator agrees with both predictions on every valid front tested: $182$ fronts over all $26$ shapes with
$2\le m\le6$, on rays inside and above the window (\texttt{verify/static\_general.py}).
A separate brute-force enumeration confirms Lemma~\ref{lem:hexsite}(b) for all $47$ shapes with
$2\le m\le8$ (\texttt{verify/static\_lattice\_order.py}).

The theorem uses the simplest perturbation. Lemma~\ref{lem:hexclear} involves only the hinges at $A$, $B$
and $C$, so on fronts with $n\ge3$ one may also turn the hinge between $U_2$ and $U_3$, and then $\alpha=0$ is
allowed: the down vectors stay exactly regular and only the up vectors carry irregularity. The channel
condition $\mu_Z>0$ is then no longer automatic. When it fails the channel between $Z$ and $E$ is closed, no
disk fits near $Q_0$, and the first disk is placed elsewhere, a triangle in all our runs
(\texttt{verify/static\_channel.py}).
